% GENERATED FILE. Edit canonical lessons, then run npm run generate:books.
% canonical-source-hash: fnv1a64:a45002d6c5e6944b
% canonical-lessons: JA-C74-maishuu, JA-C74-shuumatsu, JA-C74-owari, JA-C74-tochuu, JA-C74-han

\chapter{Every week, Weekend, End, On the way, Half past}
\label{ch:ja-every-week-weekend-end-on-the-way-half-past}
\hlchaptermodality{74}

\begin{chapteropening}
\textbf{By the end of this chapter:} \emph{I can say every week, the weekend, the end, on the way and half past.}

Complete the last lesson of chapter 74: I can say every week, the weekend, the end, on the way and half past.
\end{chapteropening}

\section[\texorpdfstring{\ja{まいしゅう} (maishuu)}{まいしゅう (maishuu)}]{\ja{まいしゅう} (maishuu) — every week}

\label{lesson:JA-C74-maishuu}



\begin{quote}
Before the new one: say the Japanese for the day before yesterday, then the Japanese for every morning.
\end{quote}

\subsection*{You'll want to know: \ja{まいしゅう}}
\textbf{\ja{まいしゅう}} — \emph{maishuu} — \textquotedblleft{}every week\textquotedblright{}.

\begin{grammarlens}[title={What you've built}]
One more word for this chapter.
\end{grammarlens}

\subsection*{Your turn}
\begin{itemize}
\raggedright
  \item \emph{Say it:} \emph{maishuu}
  \item \emph{Say it:} \emph{maishuu}, once more
  \item \emph{Say it:} say \emph{maishuu}
  \item \emph{Recall:} say the Japanese for the day before yesterday, then the Japanese for every morning, then say \emph{maishuu} again
\end{itemize}

\begin{tcolorbox}[breakable,colback=teal!4,colframe=teal!35!black,title={Before you move on}]
What does \ja{まいしゅう} mean? (\textquotedblleft{}Every week\textquotedblright{}.) Say it once more.
\end{tcolorbox}
\section[\texorpdfstring{\ja{しゅうまつ} (shuumatsu)}{しゅうまつ (shuumatsu)}]{\ja{しゅうまつ} (shuumatsu) — the weekend}

\label{lesson:JA-C74-shuumatsu}



\begin{quote}
Before the new one: say the Japanese for every morning, then the Japanese for every week.
\end{quote}

\subsection*{You'll want to know: \ja{しゅうまつ}}
\textbf{\ja{しゅうまつ}} — \emph{shuumatsu} — \textquotedblleft{}the weekend\textquotedblright{}.

\begin{grammarlens}[title={What you've built}]
One more word for this chapter.
\end{grammarlens}

\subsection*{Your turn}
\begin{itemize}
\raggedright
  \item \emph{Say it:} \emph{shuumatsu}
  \item \emph{Say it:} \emph{shuumatsu}, once more
  \item \emph{Say it:} say \emph{shuumatsu}
  \item \emph{Recall:} say the Japanese for every morning, then the Japanese for every week, then say \emph{shuumatsu} again
\end{itemize}

\begin{tcolorbox}[breakable,colback=teal!4,colframe=teal!35!black,title={Before you move on}]
What does \ja{しゅうまつ} mean? (\textquotedblleft{}The weekend\textquotedblright{}.) Say it once more.
\end{tcolorbox}
\section[\texorpdfstring{\ja{おわり} (owari)}{おわり (owari)}]{\ja{おわり} (owari) — the end}

\label{lesson:JA-C74-owari}



\begin{quote}
Before the new one: say the Japanese for every week, then the Japanese for the weekend.
\end{quote}

\subsection*{You'll want to know: \ja{おわり}}
\textbf{\ja{おわり}} — \emph{owari} — \textquotedblleft{}the end\textquotedblright{}.

\begin{grammarlens}[title={What you've built}]
One more word for this chapter.
\end{grammarlens}

\subsection*{Your turn}
\begin{itemize}
\raggedright
  \item \emph{Say it:} \emph{owari}
  \item \emph{Say it:} \emph{owari}, once more
  \item \emph{Say it:} say \emph{owari}
  \item \emph{Recall:} say the Japanese for every week, then the Japanese for the weekend, then say \emph{owari} again
\end{itemize}

\begin{tcolorbox}[breakable,colback=teal!4,colframe=teal!35!black,title={Before you move on}]
What does \ja{おわり} mean? (\textquotedblleft{}The end\textquotedblright{}.) Say it once more.
\end{tcolorbox}
\section[\texorpdfstring{\ja{とちゅう} (tochuu)}{とちゅう (tochuu)}]{\ja{とちゅう} (tochuu) — on the way, midway}

\label{lesson:JA-C74-tochuu}



\begin{quote}
Before the new one: say the Japanese for the weekend, then the Japanese for the end.
\end{quote}

\subsection*{You'll want to know: \ja{とちゅう}}
\textbf{\ja{とちゅう}} — \emph{tochuu} — \textquotedblleft{}on the way, midway\textquotedblright{}.

\begin{grammarlens}[title={What you've built}]
One more word for this chapter.
\end{grammarlens}

\subsection*{Your turn}
\begin{itemize}
\raggedright
  \item \emph{Say it:} \emph{tochuu}
  \item \emph{Say it:} \emph{tochuu}, once more
  \item \emph{Say it:} say \emph{tochuu}
  \item \emph{Recall:} say the Japanese for the weekend, then the Japanese for the end, then say \emph{tochuu} again
\end{itemize}

\begin{tcolorbox}[breakable,colback=teal!4,colframe=teal!35!black,title={Before you move on}]
What does \ja{とちゅう} mean? (\textquotedblleft{}On the way, midway\textquotedblright{}.) Say it once more.
\end{tcolorbox}
\section[\texorpdfstring{\ja{はん} (han)}{はん (han)}]{\ja{はん} (han) — half past}

\label{lesson:JA-C74-han}



\begin{quote}
Before the new one: say the Japanese for the end, then the Japanese for on the way.
\end{quote}

\subsection*{You'll want to know: \ja{はん}}
\textbf{\ja{はん}} — \emph{han} — \textquotedblleft{}half past\textquotedblright{}.

\begin{grammarlens}[title={What you've built}]
One more word for this chapter.
\end{grammarlens}

\subsection*{Your turn}
\begin{itemize}
\raggedright
  \item \emph{Say it:} \emph{han}
  \item \emph{Say it:} \emph{han}, once more
  \item \emph{Say it:} say \emph{han}
  \item \emph{Recall:} say the Japanese for the end, then the Japanese for on the way, then say \emph{han} again
\end{itemize}

\begin{tcolorbox}[breakable,colback=teal!4,colframe=teal!35!black,title={Before you move on}]
What does \ja{はん} mean? (\textquotedblleft{}Half past\textquotedblright{}.) Say it once more.
\end{tcolorbox}
